\lekcja{Sprawdziany}{Sprawdzian 1: Wielokrotności}
Imię i nazwisko: \ans[7cm]\hfill Data: \ans[3cm]\par\smallskip
{\small\color{szary}Za każde zadanie 1 punkt (razem 10 pkt). Wynik wpisz w ramkę, a obliczenia zapisz w kratkach. W zadaniach z wyborem zakreśl jedną odpowiedź.}\par
\nr{1}{\rmfamily Trzycyfrowa liczba \textbf{9\kw{}2} jest wielokrotnością liczby 26. Jaką cyfrę trzeba wpisać w okienko?}\par
\kratka{1.5}\hfill Odpowiedź: \ans[2.8cm]\par
\nr{2}{\rmfamily Oblicz \textbf{NWW(6, 5, 10)}.}\par
\kratka{1.5}\hfill Odpowiedź: \ans[2.8cm]\par
\nr{3}{\rmfamily Oblicz \textbf{NWW(6, 15)}.}\par
\kratka{1.5}\hfill Odpowiedź: \ans[2.8cm]\par
\nr{4}{\rmfamily Dziewięciokrotność liczby 14 jest dwukrotnością pewnej liczby. Jaka to liczba?}\par
\kratka{1.5}\hfill Odpowiedź: \ans[2.8cm]\par
\nr{5}{\rmfamily Liczby A, B, C to \textbf{kolejne} wielokrotności liczby 9. B = 90. Ile wynosi A?}\par
\kratka{1.5}\hfill Odpowiedź: \ans[2.8cm]\par
\nr{6}{\rmfamily Prawda czy fałsz? \\ \textit{„Każda wielokrotność liczby 4 jest wielokrotnością liczby 2.”}}\par
\smallskip \opc{A. Prawda}\quad \opc{B. Fałsz}\par
\nr{7}{\rmfamily Która z liczb jest wielokrotnością liczby \textbf{12}?}\par
\smallskip \opc{A. 160}\quad \opc{B. 161}\quad \opc{C. 171}\quad \opc{D. 156}\par
\kratka{1}
\nr{8}{\rmfamily Znajdź \textbf{największą} liczbę mniejszą od 162, która jest wielokrotnością liczby 45.}\par
\kratka{1.5}\hfill Odpowiedź: \ans[2.8cm]\par
\nr{9}{\rmfamily Od czterokrotności pewnej liczby odjęto dwukrotność tej liczby i otrzymano 4. O jakiej liczbie mowa?}\par
\kratka{1.5}\hfill Odpowiedź: \ans[2.8cm]\par
\nr{10}{\rmfamily Znajdź \textbf{najmniejszą} liczbę większą od 135, która jest wielokrotnością liczby 50.}\par
\kratka{1.5}\hfill Odpowiedź: \ans[2.8cm]\par
\par\bigskip\begin{tcolorbox}[colback=krem,colframe=pom!40,arc=6pt,fontupper=\sffamily\small]Punkty: \ans[1.5cm] / 10\hfill Ocena: \ans[3cm]\par\smallskip Skala: 0–29\% – 1,\ 30–49\% – 2,\ 50–69\% – 3,\ 70–84\% – 4,\ 85–95\% – 5,\ 96–100\% – 6.\end{tcolorbox}
\lekcja{Sprawdziany}{Sprawdzian 2: Dzielniki}
Imię i nazwisko: \ans[7cm]\hfill Data: \ans[3cm]\par\smallskip
{\small\color{szary}Za każde zadanie 1 punkt (razem 10 pkt). Wynik wpisz w ramkę, a obliczenia zapisz w kratkach. W zadaniach z wyborem zakreśl jedną odpowiedź.}\par
\nr{1}{\rmfamily Ile dzielników ma liczba \textbf{28}?}\par
\kratka{1.5}\hfill Odpowiedź: \ans[2.8cm]\par
\nr{2}{\rmfamily Znajdź najmniejszą liczbę \textbf{trzycyfrową}, której dzielnikami są liczby 12 i 5.}\par
\kratka{1.5}\hfill Odpowiedź: \ans[2.8cm]\par
\nr{3}{\rmfamily Kartkę 36 cm × 45 cm tniemy na jednakowe kwadraty o boku w pełnych centymetrach (bez odpadów). Jaki najdłuższy bok (w cm) mogą mieć kwadraty?}\par
\kratka{1.5}\hfill Odpowiedź: \ans[2.8cm]\par
\nr{4}{\rmfamily Która liczba \textbf{NIE} jest dzielnikiem liczby 30?}\par
\smallskip \opc{A. 16}\quad \opc{B. 15}\quad \opc{C. 6}\quad \opc{D. 5}\par
\kratka{1}
\nr{5}{\rmfamily Mamy 48 niebieskich kredek i 60 czerwonych kredek. Chcemy ze \textbf{wszystkich} ułożyć jednakowe zestawy. Ile najwięcej zestawów można ułożyć?}\par
\kratka{1.5}\hfill Odpowiedź: \ans[2.8cm]\par
\nr{6}{\rmfamily Oblicz sumę \textbf{wszystkich} dzielników liczby 14.}\par
\kratka{1.5}\hfill Odpowiedź: \ans[2.8cm]\par
\nr{7}{\rmfamily Oblicz \textbf{NWD(15, 35)}.}\par
\kratka{1.5}\hfill Odpowiedź: \ans[2.8cm]\par
\nr{8}{\rmfamily Prawda czy fałsz? \\ \textit{„Każdy dzielnik liczby 12 jest dzielnikiem liczby 24.”}}\par
\smallskip \opc{A. Prawda}\quad \opc{B. Fałsz}\par
\nr{9}{\rmfamily Ile jest liczb \textbf{trzycyfrowych}, których dzielnikiem jest 150?}\par
\kratka{1.5}\hfill Odpowiedź: \ans[2.8cm]\par
\nr{10}{\rmfamily \textbf{Liczba doskonała} jest równa sumie swoich dzielników mniejszych od niej. Czy 12 jest liczbą doskonałą?}\par
\smallskip \opc{A. Tak}\quad \opc{B. Nie}\par
\par\bigskip\begin{tcolorbox}[colback=krem,colframe=pom!40,arc=6pt,fontupper=\sffamily\small]Punkty: \ans[1.5cm] / 10\hfill Ocena: \ans[3cm]\par\smallskip Skala: 0–29\% – 1,\ 30–49\% – 2,\ 50–69\% – 3,\ 70–84\% – 4,\ 85–95\% – 5,\ 96–100\% – 6.\end{tcolorbox}
\lekcja{Sprawdziany}{Sprawdzian 3: Cechy podzielności}
Imię i nazwisko: \ans[7cm]\hfill Data: \ans[3cm]\par\smallskip
{\small\color{szary}Za każde zadanie 1 punkt (razem 10 pkt). Wynik wpisz w ramkę, a obliczenia zapisz w kratkach. W zadaniach z wyborem zakreśl jedną odpowiedź.}\par
\nr{1}{\rmfamily Czy rok \textbf{1850} jest przestępny?}\par
\smallskip \opc{A. Tak}\quad \opc{B. Nie}\par
\nr{2}{\rmfamily Zaginęła liczba! Jest większa od 9400 i mniejsza od 9500, dzieli się przez 9 i przez 10. Jaka to liczba?}\par
\kratka{1.5}\hfill Odpowiedź: \ans[2.8cm]\par
\nr{3}{\rmfamily Która liczba jest podzielna przez \textbf{15}?}\par
\smallskip \opc{A. 4200}\quad \opc{B. 8103}\quad \opc{C. 3823}\quad \opc{D. 1898}\par
\kratka{1}
\nr{4}{\rmfamily Nie wykonując dzielenia, podaj resztę z dzielenia liczby \textbf{26\,646} przez 5.}\par
\kratka{1.5}\hfill Odpowiedź: \ans[2.8cm]\par
\nr{5}{\rmfamily Prawda czy fałsz? \\ \textit{„Każda liczba podzielna przez 3 jest nieparzysta.”}}\par
\smallskip \opc{A. Prawda}\quad \opc{B. Fałsz}\par
\nr{6}{\rmfamily Jaką cyfrę trzeba wpisać w okienko, żeby liczba \textbf{9\kw{}04} była podzielna przez 9?}\par
\kratka{1.5}\hfill Odpowiedź: \ans[2.8cm]\par
\nr{7}{\rmfamily Suma cyfr pewnej liczby wynosi \textbf{2160}. Czy ta liczba jest podzielna przez 9?}\par
\smallskip \opc{A. Tak}\quad \opc{B. Nie}\par
\nr{8}{\rmfamily Jaką \textbf{najmniejszą} cyfrę można wpisać w okienko, żeby liczba \textbf{5\kw{}3} była podzielna przez 3?}\par
\kratka{1.5}\hfill Odpowiedź: \ans[2.8cm]\par
\nr{9}{\rmfamily Która liczba jest podzielna przez \textbf{10}?}\par
\smallskip \opc{A. 25\,500}\quad \opc{B. 25\,487}\quad \opc{C. 74\,408}\quad \opc{D. 98\,618}\par
\kratka{1}
\nr{10}{\rmfamily Która liczba jest podzielna przez \textbf{6}?}\par
\smallskip \opc{A. 5794}\quad \opc{B. 1306}\quad \opc{C. 2526}\quad \opc{D. 2158}\par
\kratka{1}
\par\bigskip\begin{tcolorbox}[colback=krem,colframe=pom!40,arc=6pt,fontupper=\sffamily\small]Punkty: \ans[1.5cm] / 10\hfill Ocena: \ans[3cm]\par\smallskip Skala: 0–29\% – 1,\ 30–49\% – 2,\ 50–69\% – 3,\ 70–84\% – 4,\ 85–95\% – 5,\ 96–100\% – 6.\end{tcolorbox}
\lekcja{Sprawdziany}{Sprawdzian 4: Liczby pierwsze}
Imię i nazwisko: \ans[7cm]\hfill Data: \ans[3cm]\par\smallskip
{\small\color{szary}Za każde zadanie 1 punkt (razem 10 pkt). Wynik wpisz w ramkę, a obliczenia zapisz w kratkach. W zadaniach z wyborem zakreśl jedną odpowiedź.}\par
\nr{1}{\rmfamily Którą cyfrą jedności \textbf{nie może} kończyć się liczba pierwsza większa od 10?}\par
\smallskip \opc{A. 7}\quad \opc{B. 1}\quad \opc{C. 6}\quad \opc{D. 3}\par
\kratka{1}
\nr{2}{\rmfamily Prawda czy fałsz? \\ \textit{„Liczba 2 jest jedyną parzystą liczbą pierwszą.”}}\par
\smallskip \opc{A. Prawda}\quad \opc{B. Fałsz}\par
\nr{3}{\rmfamily Podaj dowolny dzielnik liczby \textbf{221} różny od 1 i od 221 (to pokazuje, że jest złożona).}\par
\kratka{1.5}\hfill Odpowiedź: \ans[2.8cm]\par
\nr{4}{\rmfamily Jaka jest najmniejsza liczba pierwsza \textbf{większa} od 82?}\par
\kratka{1.5}\hfill Odpowiedź: \ans[2.8cm]\par
\nr{5}{\rmfamily Liczba 59 jest pierwsza. Oblicz \textbf{NWD(61, 59)}.}\par
\kratka{1.5}\hfill Odpowiedź: \ans[2.8cm]\par
\nr{6}{\rmfamily Która z liczb jest \textbf{pierwsza}?}\par
\smallskip \opc{A. 69}\quad \opc{B. 117}\quad \opc{C. 83}\quad \opc{D. 147}\par
\kratka{1}
\nr{7}{\rmfamily Ile jest liczb pierwszych większych od 30 i mniejszych od 60?}\par
\kratka{1.5}\hfill Odpowiedź: \ans[2.8cm]\par
\nr{8}{\rmfamily Która z liczb jest \textbf{złożona}?}\par
\smallskip \opc{A. 113}\quad \opc{B. 101}\quad \opc{C. 37}\quad \opc{D. 121}\par
\kratka{1}
\nr{9}{\rmfamily Ile dzielników ma liczba \textbf{49} (7 · 7)?}\par
\kratka{1.5}\hfill Odpowiedź: \ans[2.8cm]\par
\nr{10}{\rmfamily Ile jest liczb pierwszych większych od 40 i mniejszych od 60?}\par
\kratka{1.5}\hfill Odpowiedź: \ans[2.8cm]\par
\par\bigskip\begin{tcolorbox}[colback=krem,colframe=pom!40,arc=6pt,fontupper=\sffamily\small]Punkty: \ans[1.5cm] / 10\hfill Ocena: \ans[3cm]\par\smallskip Skala: 0–29\% – 1,\ 30–49\% – 2,\ 50–69\% – 3,\ 70–84\% – 4,\ 85–95\% – 5,\ 96–100\% – 6.\end{tcolorbox}
\lekcja{Sprawdziany}{Sprawdzian 5: Rozkład na czynniki}
Imię i nazwisko: \ans[7cm]\hfill Data: \ans[3cm]\par\smallskip
{\small\color{szary}Za każde zadanie 1 punkt (razem 8 pkt). Wynik wpisz w ramkę, a obliczenia zapisz w kratkach. W zadaniach z wyborem zakreśl jedną odpowiedź.}\par
\nr{1}{\rmfamily Jaka liczba ma rozkład na czynniki pierwsze \textbf{2 · 2 · 5}?}\par
\kratka{1.5}\hfill Odpowiedź: \ans[2.8cm]\par
\nr{2}{\rmfamily Jaki jest \textbf{największy} czynnik pierwszy liczby 384?}\par
\kratka{1.5}\hfill Odpowiedź: \ans[2.8cm]\par
\nr{3}{\rmfamily Uzupełnij rozkład: \textbf{45 = 3 · \kw{} · 5}}\par
\kratka{1.5}\hfill Odpowiedź: \ans[2.8cm]\par
\nr{4}{\rmfamily Ile czynników (licząc powtórzenia) ma rozkład liczby \textbf{298} na czynniki pierwsze?}\par
\kratka{1.5}\hfill Odpowiedź: \ans[2.8cm]\par
\nr{5}{\rmfamily 20 = 2 · 2 · 5, ~ 9 = 3 · 3. Oblicz \textbf{NWW(20, 9)}.}\par
\kratka{1.5}\hfill Odpowiedź: \ans[2.8cm]\par
\nr{6}{\rmfamily 8 = 2 · 2 · 2, ~ 20 = 2 · 2 · 5. Oblicz \textbf{NWD(8, 20)}.}\par
\kratka{1.5}\hfill Odpowiedź: \ans[2.8cm]\par
\nr{7}{\rmfamily 18 = 2 · 3 · 3, ~ 15 = 3 · 5. Oblicz \textbf{NWW(18, 15)}.}\par
\kratka{1.5}\hfill Odpowiedź: \ans[2.8cm]\par
\nr{8}{\rmfamily 42 = 2 · 3 · 7, ~ 18 = 2 · 3 · 3. Oblicz \textbf{NWD(42, 18)}.}\par
\kratka{1.5}\hfill Odpowiedź: \ans[2.8cm]\par
\par\bigskip\begin{tcolorbox}[colback=krem,colframe=pom!40,arc=6pt,fontupper=\sffamily\small]Punkty: \ans[1.5cm] / 8\hfill Ocena: \ans[3cm]\par\smallskip Skala: 0–29\% – 1,\ 30–49\% – 2,\ 50–69\% – 3,\ 70–84\% – 4,\ 85–95\% – 5,\ 96–100\% – 6.\end{tcolorbox}
\lekcja{Sprawdziany}{Sprawdzian zbiorczy}
Imię i nazwisko: \ans[7cm]\hfill Data: \ans[3cm]\par\smallskip
{\small\color{szary}Za każde zadanie 1 punkt (razem 16 pkt). Wynik wpisz w ramkę, a obliczenia zapisz w kratkach. W zadaniach z wyborem zakreśl jedną odpowiedź.}\par
\nr{1}{\rmfamily Liczba 23 jest pierwsza. Oblicz \textbf{NWD(191, 23)}.}\par
\kratka{1.5}\hfill Odpowiedź: \ans[2.8cm]\par
\nr{2}{\rmfamily Nie wykonując dzielenia, podaj resztę z dzielenia liczby \textbf{27\,918} przez 9.}\par
\kratka{1.5}\hfill Odpowiedź: \ans[2.8cm]\par
\nr{3}{\rmfamily Suma cyfr pewnej liczby wynosi \textbf{9186}. Czy ta liczba jest podzielna przez 3?}\par
\smallskip \opc{A. Tak}\quad \opc{B. Nie}\par
\nr{4}{\rmfamily Oblicz \textbf{NWW(5, 8, 10)}.}\par
\kratka{1.5}\hfill Odpowiedź: \ans[2.8cm]\par
\nr{5}{\rmfamily Mamy 18 cukierków i 24 lizaków. Chcemy ze \textbf{wszystkich} ułożyć jednakowe zestawy. Ile najwięcej zestawów można ułożyć?}\par
\kratka{1.5}\hfill Odpowiedź: \ans[2.8cm]\par
\nr{6}{\rmfamily Znajdź \textbf{największą} liczbę mniejszą od 632, która jest wielokrotnością liczby 80.}\par
\kratka{1.5}\hfill Odpowiedź: \ans[2.8cm]\par
\nr{7}{\rmfamily Która liczba \textbf{NIE} jest dzielnikiem liczby 40?}\par
\smallskip \opc{A. 20}\quad \opc{B. 10}\quad \opc{C. 2}\quad \opc{D. 22}\par
\kratka{1}
\nr{8}{\rmfamily Która liczba jest podzielna przez \textbf{2}?}\par
\smallskip \opc{A. 83\,067}\quad \opc{B. 56\,537}\quad \opc{C. 41\,775}\quad \opc{D. 42\,446}\par
\kratka{1}
\nr{9}{\rmfamily Prawda czy fałsz? \\ \textit{„Liczba 13 ma dokładnie dwa dzielniki.”}}\par
\smallskip \opc{A. Prawda}\quad \opc{B. Fałsz}\par
\nr{10}{\rmfamily Uzupełnij rozkład: \textbf{56 = \kw{} · 2 · 2 · 7}}\par
\kratka{1.5}\hfill Odpowiedź: \ans[2.8cm]\par
\nr{11}{\rmfamily Od pięciokrotności pewnej liczby odjęto trzykrotność tej liczby i otrzymano 26. O jakiej liczbie mowa?}\par
\kratka{1.5}\hfill Odpowiedź: \ans[2.8cm]\par
\nr{12}{\rmfamily Jaka jest najmniejsza liczba pierwsza \textbf{większa} od 39?}\par
\kratka{1.5}\hfill Odpowiedź: \ans[2.8cm]\par
\nr{13}{\rmfamily \textbf{Liczba doskonała} jest równa sumie swoich dzielników mniejszych od niej. Czy 18 jest liczbą doskonałą?}\par
\smallskip \opc{A. Tak}\quad \opc{B. Nie}\par
\nr{14}{\rmfamily Ile czynników (licząc powtórzenia) ma rozkład liczby \textbf{228} na czynniki pierwsze?}\par
\kratka{1.5}\hfill Odpowiedź: \ans[2.8cm]\par
\nr{15}{\rmfamily Znajdź \textbf{najmniejszą} liczbę większą od 381, która jest wielokrotnością liczby 40.}\par
\kratka{1.5}\hfill Odpowiedź: \ans[2.8cm]\par
\nr{16}{\rmfamily Która z liczb jest \textbf{złożona}?}\par
\smallskip \opc{A. 109}\quad \opc{B. 59}\quad \opc{C. 169}\quad \opc{D. 29}\par
\kratka{1}
\par\bigskip\begin{tcolorbox}[colback=krem,colframe=pom!40,arc=6pt,fontupper=\sffamily\small]Punkty: \ans[1.5cm] / 16\hfill Ocena: \ans[3cm]\par\smallskip Skala: 0–29\% – 1,\ 30–49\% – 2,\ 50–69\% – 3,\ 70–84\% – 4,\ 85–95\% – 5,\ 96–100\% – 6.\end{tcolorbox}